\chapter{Host Organisation and MFJA Robotics Platforms}
\label{chap:host}

\section{Université de Toulouse}

The internship was formally hosted by Université de Toulouse and carried out at
the MFJA site, 1 rue Tarfaya, Toulouse. As a public university, the institution
combines higher education, research and knowledge transfer. Its official
Nomenclature d'activités française/Activité principale exercée (NAF/APE) code
is 85.42Z \citep{univtoulousesirene}, which INSEE classifies as \emph{higher
education}, including first-, second- and third-cycle university programmes
\citep{insee8542z}.

Université de Toulouse has operated as an experimental public institution since
January 2025. Official figures published for 2026 indicate approximately 39,000
students and 4,200 staff \citep{univtoulousefigures}. These figures concern the
university as a whole, while the internship took place within the smaller MFJA
environment.

\section{Maison de la Formation Jacqueline Auriol}

The MFJA is a shared centre for engineering education on the Toulouse Aerospace
innovation campus. Its activities cover research and innovation, knowledge and
skills transfer, and initial and continuing education. According to its official
presentation, the site brings together more than 1,500 students from 15
mechanical-engineering programmes associated with Université de Toulouse, INSA
and ISAE. It supports work in mechanical engineering, manufacturing,
aeronautics, automation and robotics in collaboration with industrial partners
and research laboratories \citep{mfjaofficial}.

Because the facilities are shared across teaching and research activities, the
software developed for the platform must support several configurations, remain
usable by different operators and be maintainable beyond a single development
environment. These requirements influenced the package and launch architecture
presented in Chapters~5--6.

\section{The Room 315 robotics environment}

Room~315 contains two mirrored flexible transport cells, each with four work
positions associated with industrial robots. In the configuration used for
this project, the right rail serves Yaskawa HC10DT positions~1--2 and Stäubli
TX2 positions~3--4, while the left rail serves Yaskawa HC10 positions~1--2 and
KUKA KR6 positions~3--4. Each circuit contains outer and inner branches, four
switching zones, indexing detectors and stoppers. A shuttle may be empty or
carry a payload.

This arrangement brings together several engineering disciplines:

\begin{itemize}
  \item mechanical layout, computer-aided design (CAD), kinematics and
  collision geometry;
  \item industrial automation, sensors, stoppers and route switching;
  \item robot middleware, namespacing, timing and distributed interfaces;
  \item planning, scheduling and multi-agent resource conflicts;
  \item computer vision, data collection and model evaluation;
  \item human--robot interaction through task-level language.
\end{itemize}

Together, these disciplines make Room~315 a practical platform for studying the
integration of robotics, transport, perception and control.
\Cref{fig:room315-orientation} shows the simulated layout; Chapters~7--8 define
the rail topology, device semantics and supervisory constraints.

\begin{figure}[htbp]
  \centering
  \includegraphics[width=0.96\textwidth]{f04_room315_gazebo_overview.png}
  \caption[Gazebo overview of Room~315 and its four robot workstations.]
  {Gazebo view of Room~315. Four industrial-robot workstations surround two
  mirrored rail-transport cells containing eight stopped shuttles. The measured
  rail geometry and device locations are presented in \cref{chap:transport}.}
  \label{fig:room315-orientation}
\end{figure}

\section{Organisation and responsibilities}

\Cref{fig:organisation} presents the supervision and responsibilities directly
associated with the internship.

\begin{figure}[htbp]
\centering
\resizebox{0.94\textwidth}{!}{%
\begin{tikzpicture}[
  node distance=9mm and 16mm,
  box/.style={draw,rounded corners,align=center,minimum height=10mm,text width=42mm,fill=mfjalight},
  role/.style={draw,rounded corners,align=left,minimum height=13mm,text width=58mm,fill=white},
  arr/.style={-{Latex[length=2.2mm]},thick,mfjablue}
]
  \node[box] (host) {Université de Toulouse\\MFJA};
  \node[box,below=of host] (tutor) {Host supervisors\\Joint scientific supervision};
  \node[box,below=of tutor] (student) {Ali Ibrahim\\M2 intern};
  \node[role,below left=of student] (ops) {Operational knowledge\\Platform access\\Review of platform constraints};
  \node[role,below right=of student] (dev) {Architecture and implementation\\Testing and experiments\\Documentation and handover};
  \draw[arr] (host) -- (tutor);
  \draw[arr] (tutor) -- (student);
  \draw[arr] (student) -- (ops);
  \draw[arr] (student) -- (dev);
\end{tikzpicture}
}
\caption{Organisation of the internship and distribution of responsibilities.}
\label{fig:organisation}
\end{figure}

The host supervisors provided joint scientific direction and reviewed choices
related to the platform. MFJA personnel supported access, scheduling and
practical organisation. I was responsible for the architecture, implementation,
evaluation, documentation and handover described in \cref{chap:mission}.

\section{Position of the mission within the host}

The mission connects platform operation with research prototyping. Platform
managers, students and researchers can use the system to launch robots and rail
scenarios reproducibly. The same foundation can support synthetic-data
generation, task-planning experiments and comparisons between simulation
configurations. Possible uses include a teaching exercise on ROS~2 topics, a
planner benchmark, a visual-perception dataset or a simulation demonstrator.

This range of uses motivated the emphasis on reusable configuration, explicit
interfaces and documented operating procedures.

\section{Economic and strategic value}

Because MFJA is a public teaching and research platform, the mission's value is
assessed through practical reuse, training capacity and reduced integration
effort. The main benefits are:

\begin{itemize}
  \item \textbf{reuse}: one model and launch framework avoids recreating an
  integration for each practical class or research project;
  \item \textbf{risk reduction}: simulation and supervised execution check
  naming, topology and commands before access to costly equipment;
  \item \textbf{asset visibility}: a digital twin makes robot capabilities
  available for remote preparation, including capabilities that might otherwise
  remain under-used;
  \item \textbf{training capacity}: repeatable scenario generation creates
  data and exercises without monopolising the physical cell;
  \item \textbf{maintenance}: typed interfaces, tests and documentation reduce
  dependence on knowledge held by a single developer;
  \item \textbf{technology transfer}: the MFJA can demonstrate current ROS~2,
  planning and AI methods while maintaining clear safety boundaries.
\end{itemize}

These benefits align with the MFJA's institutional role of sharing costly
engineering resources and connecting education with contemporary industrial
methods.
